\documentclass[lettersize,journal]{IEEEtran}


\IEEEoverridecommandlockouts
% The preceding line is only needed to identify funding in the first footnote. If that is unneeded, please comment it out.
\usepackage{cite}
\usepackage{amsmath,amssymb,amsfonts}
\usepackage{algorithmic}
\usepackage{graphicx}
\usepackage{textcomp}
\usepackage{xcolor}
% Color codes for reviewer-specific revisions
\definecolor{AEcolor}{RGB}{0,0,200}        % Associate Editor: blue
\definecolor{R1color}{RGB}{200,0,0}        % Reviewer 1: red
\definecolor{R2color}{RGB}{180,100,0}      % Reviewer 2: orange-brown
\definecolor{R3color}{RGB}{0,140,80}       % Reviewer 3: teal-green
\usepackage{booktabs, longtable}
\usepackage{enumitem,threeparttable}
\usepackage[most]{tcolorbox}
\tcbuselibrary{listings, breakable}
\usepackage{stfloats}
\usepackage{multirow}
\usepackage{newtxtext}
\usepackage{hyperref}
\usepackage{subcaption}
\usepackage{float} 
\usepackage{caption}
\usepackage{array}
\usepackage{arydshln} % for \hdashline in tables
\usepackage{etoolbox} 
\usepackage{placeins}

\def\BibTeX{{\rm B\kern-.05em{\sc i\kern-.025em b}\kern-.08em
    T\kern-.1667em\lower.7ex\hbox{E}\kern-.125emX}}
\begin{document}

\title{DK-Root: A Joint Data-and-Knowledge-Driven Framework for Root Cause Analysis of QoE Degradations in Mobile Networks}

\author{Qizhe Li, Haolong Chen, Jiansheng Li, Shuqi Chai, \\
Xuan Li, Yuzhou Hou, Xinhua Shao, Fangfang Li, Kaifeng Han, Guangxu Zhu,~\IEEEmembership{Member,~IEEE}%
\thanks{Qizhe Li, Haolong Chen, Jiansheng Li, Shuqi Chai and Guangxu Zhu are with Shenzhen Research Institute of Big Data, Shenzhen, China, and the School of Science and Engineering, The Chinese University of Hong Kong, Shenzhen, China (Emails: \{qizheli, haolongchen1, jianshengli\}@link.cuhk.edu.cn, \{schai, gxzhu\}@sribd.cn).}%
\thanks{Xuan Li and Yuzhou Hou are with the Experience Lab, Huawei Technologies Co., Ltd., Shenzhen, China (Emails: \{lixuan69, houyuzhou\}@huawei.com).}%
\thanks{Xinhua Shao and Fangfang Li are with China Telecommunications Group Co., Ltd., Beijing, China (Email: \{shaoxh, liff\}@chinatelecom.cn).}%
\thanks{Kaifeng Han is with the China Academy of Information and Communications Technology, Beijing, China (Email: hankaifeng@caict.ac.cn).}%
\thanks{Corresponding author: Guangxu Zhu (gxzhu@sribd.cn).}%
}

% The paper headers
\markboth{Journal of \LaTeX\ Class Files,~Vol.~14, No.~8, August~2021}%
{Shell \MakeLowercase{\textit{et al.}}: A Sample Article Using IEEEtran.cls for IEEE Journals}

\IEEEpubid{0000--0000/00\$00.00~\copyright~2021 IEEE}
% Remember, if you use this you must call \IEEEpubidadjcol in the second
% column for its text to clear the IEEEpubid mark.

\maketitle


\begin{abstract}
Diagnosing the root causes of Quality of Experience (QoE) degradations in operational mobile networks is challenging due to complex cross-layer interactions among kernel performance indicators (KPIs) and the scarcity of reliable expert annotations. Although rule-based heuristics can generate labels at scale, they are noisy and coarse-grained, limiting the accuracy of purely data-driven approaches. To address this, we propose DK-Root, a joint data–and-knowledge-driven framework that unifies scalable weak supervision with precise expert guidance for root-cause analysis. DK-Root first pretrains an encoder via contrastive representation learning using abundant rule-based labels while explicitly denoising their noise through a supervised contrastive objective. To supply task-faithful data augmentation, we introduce a class-conditional diffusion model that generates KPIs sequences preserving root-cause semantics, and by controlling reverse diffusion steps, it produces weak and strong augmentations that improve intra-class compactness and inter-class separability. Finally, the encoder and the lightweight classifier are jointly fine-tuned with scarce expert-verified labels to sharpen decision boundaries. 
Experiments on a real-world operator-grade dataset show that DK-Root consistently outperforms strong competing methods, delivers more balanced and discriminative class-level predictions for confusable root causes, and approaches comparable root-cause analysis performance with approximately one-third to one-half fewer expert annotations.

\end{abstract}


\begin{IEEEkeywords}
Mobile communication, root cause analysis, deep learning, time series.
\end{IEEEkeywords}


\section{Introduction}
\IEEEPARstart{E}{nsuring} high Quality of Experience (QoE) in mobile networks is increasingly difficult as architectures grow more complex, with tightly coupled protocol layers and heterogeneous elements. These cross-layer dependencies obscure how degradations arise and propagate, slowing diagnosis and remediation. At the same time, 6G aims for autonomous operation and embedded intelligence, which further increases the need for reliable, data-driven root-cause analysis. Recent directions such as statistical digital twins~\cite{luo2023srcon} seek closed-loop integration between physical infrastructure and virtual models for real-time inference and adaptation. Within this context, accurately identifying the root causes of QoE drops (e.g., excessive latency or low throughput) from multivariate KPIs time series is pivotal. As illustrated in Figure~\ref{fig:background}, degradations can stem from diverse factors including interference, weak coverage, or resource congestion, and their signatures often span the physical (PHY), medium access control (MAC), radio link control (RLC), and packet data convergence protocol (PDCP) layers. However, constructing robust diagnostic models remains challenging due to scarce, costly expert annotations and the brittleness of rule-based heuristics, which scale labeling but introduce noise and coarse granularity. These realities call for frameworks that can leverage abundant weak supervision without sacrificing semantic fidelity, while selectively incorporating expert knowledge to refine decision boundaries.


\begin{figure}
    \centering
    \includegraphics[width=\linewidth]{Figure/background.pdf}
    \caption{Examples of poor communication QoE caused by different root causes.}
    \label{fig:background}
\end{figure}


\IEEEpubidadjcol

Existing methods for QoE root-cause analysis broadly fall into two categories: rule-based and learning-based approaches. Rule-based methods encode domain knowledge into handcrafted logic, such as threshold conditions on signal-to-interference-plus-noise ratio (SINR) or reference signal received power (RSRP), and are widely used due to their simplicity~\cite{ciocarlie2014feasibility, moulay2020novel, yuan2020anomaly}. However, they tend to be brittle when facing subtle, compound, or previously unseen degradation patterns. In contrast, deep learning methods have demonstrated the ability to extract complex cross-layer features from KPIs streams~\cite{chawla2020interpretable, sun2024spotlight}. Yet, their performance strongly depends on the availability of high-quality annotated training data, which are difficult to obtain in real-world network environments.

A key bottleneck in mobile network diagnosis lies in the scarcity of reliable root-cause annotations. Each communication record comprises multivariate KPIs sequences spanning multiple protocol layers (e.g., PHY, MAC, RLC, PDCP), whose semantics are neither intuitive nor visually interpretable. Unlike image or audio data, understanding these KPIs sequences, especially their temporal and cross-layer interactions, requires deep domain expertise in mobile protocol behaviors and fault propagation \cite{chen2024overview}. As a result, expert-verified root-cause labels are both expensive and limited in scale.
To scale up supervision, practitioners often resort to rule-based heuristics derived from KPIs thresholds and logical expressions to generate large quantities of weak labels. However, these rule-based labels tend to be coarse-grained and noisy due to their simplistic assumptions and lack of context awareness. This leads to a significant label quality gap between abundant but imprecise rule labels and accurate yet scarce expert annotations.

To bridge this gap, semi-supervised learning (SSL) has gained increasing traction. In particular, semi-supervised contrastive learning (SSCL) has emerged as a powerful paradigm for time-series tasks~\cite{zhang2024self}, as it enables discriminative representation learning using a small number of high-quality labels. By encouraging samples with similar degradation patterns to cluster together (positive pairs) while pushing dissimilar ones apart (negative pairs), SSCL facilitates robust feature learning under weak supervision.

However, the effectiveness of SSCL heavily relies on the design of data augmentations. In the time-series domain, augmentation-based contrastive learning is commonly used~\cite{wen2020time, peng2021fault, luo2023time}, with transformations such as jittering, scaling, or time-warping applied to generate diverse views of the same sequence. Yet, these generic augmentations are often task-agnostic and risk violating the physical or protocol-layer constraints embedded in KPIs. For example, injecting noise into RSRP or shuffling SINR sequences may disrupt signal continuity or scheduling semantics, leading to semantic drift and degraded diagnostic performance.

These challenges motivate the need for communication-aware and cause-consistent data augmentation strategies that preserve the structural integrity and semantic patterns of KPIs sequences. In particular, generative models that can synthesize augmented views conditioned on specific degradation types offer a promising direction, as they enable the creation of diverse yet semantically aligned samples for contrastive training.


To address the challenge of root-cause analysis for QoE degradations under limited supervision, we propose \textbf{DK-Root}, a joint data-and-knowledge-driven framework tailored for this task in mobile networks. DK-Root is motivated by the observation that large-scale rule-based root-cause annotations represent a scalable but noisy form of data-driven supervision, whereas expert-verified diagnostic results provide sparse but reliable knowledge-driven supervision. To effectively leverage both sources, we unify coarse-grained representation learning and fine-grained expert refinement into a single framework. Our key contributions are as follows:


\begin{itemize}
    \item \textbf{A Joint Data-and-Knowledge QoE Root-Cause Analysis Framework:} We design DK-Root to unify two complementary supervision sources: (i) rule-based labels that enable scalable but coarse data-driven supervision, and (ii) expert-verified labels that provide precise but scarce knowledge-driven guidance. To bridge the gap between data abundance and knowledge scarcity, we decouple feature representation learning and final classification, allowing the model to benefit from both broad heuristics and high-quality human expertise.

    \item \textbf{Conditional Generative Augmentation under Expert Guidance:} We introduce a conditional diffusion-based generative module that synthesizes class-consistent KPIs sequences. Expert labels are first embedded and fused with KPIs representations through a fusion network, forming semantically aligned conditional inputs. We further control the reverse diffusion time step to produce augmented samples of varying semantic intensities: weak augmentation (small time steps) ensures stability, while strong augmentation (large time steps) introduces higher diversity. These expert-conditioned augmentations improve intra-class compactness and enhance generalization under limited supervision.

    \item \textbf{Data-Aware Contrastive Representation Learning with Knowledge-Based Refinement:} We employ a supervised contrastive learning strategy to construct a degradation-aware embedding space. By pulling together samples with similar rule-based labels and pushing apart those with different root causes, the model learns to denoise noisy supervision while capturing the underlying structure of degradation patterns. This contrastive pretraining is followed by expert-guided fine-tuning using a small number of high-fidelity labels, which further refines the classification boundary and improves root-cause discrimination in noisy and low-resource settings.



    \item \textbf{Comprehensive Validation in Realistic Settings:} We evaluate DK-Root on a real-world mobile network dataset annotated with both rule-based and expert labels. Through extensive experiments, including baseline comparisons, sensitivity analyses on expert-label and rule-label quantities, controlled rule-noise studies, semantic evaluation of generated samples, class-level ablation analyses, and deployment-cost measurements, we show that DK-Root improves root-cause identification performance, reduces expert-label requirements, mitigates confusion among similar causes, and remains practical for deployment.
\end{itemize}


The remainder of this paper is organized as follows. Section~\ref{sec:related_work} reviews related work. Section~\ref{task description} defines the problem and notation. Section~\ref{sec:proposed method} presents our proposed framework. Section~\ref{sec:experiments} reports experimental results and analysis. Section~\ref{conclusion} concludes the paper and discusses future work.



\section{Related Work}
\label{sec:related_work}


\subsection{QoE Anomaly Detection}

QoE root-cause analysis has emerged as a key research topic in mobile communication networks, driven by the growing need for automated and intelligent network operation. Existing approaches can be broadly divided into two categories: traditional rule-based methods and modern data-driven learning models.

\noindent\textbf{Traditional Rule-Based Methods.}
Early works focused on encoding expert knowledge into static rule sets or statistical models for anomaly detection and root-cause inference. Ciocarlie et al.~\cite{ciocarlie2014feasibility} examined the feasibility of deploying anomaly detection systems within operational cellular networks, proposing probabilistic rule-based frameworks that rely on Markov logic networks. Their approach incorporated domain heuristics to detect network malfunctions, but required manual rule design and struggled with generalization across cell types and deployment contexts. Moulay et al.~\cite{moulay2020novel} proposed a decision-tree-based system for automated detection and classification of networking anomalies. While computationally efficient, their method is sensitive to pre-defined thresholding logic and suffers from rigidity in the face of unseen faults or cross-layer interactions. Yuan et al.~\cite{yuan2020anomaly} explored the integration of AI into the root-cause analysis workflow. They proposed a hybrid anomaly detection system combining statistical profiling and shallow machine learning, with a focus on deployment viability in large-scale radio networks. However, their approach still heavily relies on manually defined KPIs patterns, limiting scalability. Ramírez et al.~\cite{ramirez2023explainable} contributed to the explainability of anomaly detection by incorporating interpretable machine learning models (e.g., decision trees and SHAP values) for mobile network performance monitoring. While improving transparency, the method inherits the generalization limitations of rule-centric systems and cannot fully leverage unlabeled data.

\noindent\textbf{Deep Learning and Unsupervised Models.}
More recently, researchers have turned to deep learning to exploit the rich structure in time-series KPIs data. Chawla et al.~\cite{chawla2020interpretable} proposed an interpretable unsupervised anomaly detection method for RAN cell traces, using autoencoder-based architectures to uncover latent anomalies in network traces. Their framework emphasized explainability through feature attribution, but its unsupervised nature limits precise root-cause classification. Sun et al.~\cite{sun2024spotlight} developed SpotLight, a high-accuracy, explainable, and efficient anomaly detection system tailored for Open RAN environments. The model leverages hierarchical spatial-temporal feature extraction combined with outlier scoring to detect performance drops at scale. However, although SpotLight offers strong anomaly detection capabilities, it stops short of fine-grained root-cause categorization, leaving the diagnosis problem open.

Despite recent advances, existing approaches face a practical trade-off: rule-based methods offer interpretability but limited flexibility, whereas deep learning models usually require more reliable labels than are available in operational settings. In addition, prior studies generally do not discuss how augmentation should preserve the communication semantics of KPIs such as RSRP and SINR. Our work focuses on this setting by combining expert-labeled and rule-labeled data within a semi-supervised contrastive framework and by using class-conditional augmentation to better preserve root-cause-related KPI patterns. This design aims to improve label efficiency while maintaining diagnostic fidelity in mobile-network root-cause analysis.



\subsection{Time-series Classification}

\noindent\textbf{Traditional Methods.}  
Classic TSC methods often rely on distance-based comparisons such as 1-Nearest Neighbor with Dynamic Time Warping (DTW), which remain strong baselines~\cite{bagnall2017great}. Enhancements like shapelets or symbolic representations can boost performance without extensive learning~\cite{baydogan2015learning}. These methods offer interpretability and simplicity, but struggle when handling high-dimensional, multimodal sequences and lack the ability to learn temporal abstractions.

\noindent\textbf{Deep Learning for TSC.}  
Deep learning has broadened the modeling choices for TSC. CNNs, recurrent models such as LSTM/GRU, and more recent Transformer-based architectures have all shown competitive performance on benchmark datasets~\cite{ismail2019deep, ismail2020inceptiontime, fawaz2020deep}. These methods can learn hierarchical temporal representations, but their effectiveness often depends on sufficient labeled data and inductive biases that match the target signals. This limitation has motivated semi-supervised and self-supervised TSC research, where unlabeled sequences are used to improve representations when annotations are scarce. Contrastive learning frameworks such as TS-TCC (Temporal and Contextual Contrast) and its extension CA-TCC (Class-Aware TS-TCC) enforce consistency across augmented views of the same sequence~\cite{eldele2021time, eldele2023self}. TS2Vec further introduces hierarchical contrastive objectives to learn embeddings at multiple temporal scales~\cite{yue2022ts2vec}. However, these methods typically rely on generic augmentations such as jittering, scaling, or warping.
Beyond contrastive learning, low-label TSC has also been studied through pseudo-labeling, self-training, reconstruction, and forecasting objectives, which can improve robustness when paired with suitable architectures and augmentation strategies~\cite{eldele2023self}. 

Despite these advances, many semi-supervised TSC methods remain domain-agnostic and overlook signal-specific constraints. For network KPIs, augmentations that distort amplitude or time structure can corrupt physically meaningful patterns and may weaken downstream root-cause identification.

% {\color{R3color}\noindent\textbf{Learning with Noisy Labels.}
% When supervision is derived from heuristic rules rather than human experts, the resulting labels can be substantially noisy. The noisy-label learning literature has proposed several strategies to mitigate this issue. \textit{Co-teaching}~\cite{han2018coteaching} trains two networks in parallel, each selecting small-loss samples to teach the other, effectively filtering out noisy examples through mutual denoising. \textit{DivideMix}~\cite{li2020dividemix} combines a Gaussian Mixture Model for dynamically separating clean and noisy samples with MixMatch-style semi-supervised learning on the separated data. These methods have shown effectiveness on image classification benchmarks; however, their applicability to multivariate time-series data with protocol-layer constraints remains underexplored. In contrast, our approach handles label noise at the representation level: by using noisy rule labels to guide contrastive learning rather than cross-entropy classification, the embedding space naturally becomes more tolerant of individual label errors.}

Building on these insights, we propose a domain-aware semi-supervised contrastive learning framework tailored for QoE root-cause analysis. By integrating expert knowledge into augmentation design and contrastive objectives, we ensure that learned embeddings maintain the rich temporal and protocol-layer semantics of network KPIs sequences. Moreover, our unified learning pipeline fuses expert-labeled and rule-labeled samples, delivering a balance of label efficiency and precision.



\section{Problem Formulation}
\label{task description}

\begin{figure*}[!t]
    \centering
    \includegraphics[width=1\linewidth]{Figure/Bigpicture.png}
    \caption{Overview of QoE root-cause analysis in mobile networks.}
    \label{fig:bigpicture}
\end{figure*}

We formalize the root-cause analysis of QoE degradations in mobile networks as a classification learning problem over multi-dimensional KPIs time-series data, where each sample corresponds to a sequence of communication performance metrics collected across multiple protocol layers. 
In this paper, we focus on a representative case of high-latency scenarios, {\color{R1color}which is one of the most critical QoE degradations in 5G and beyond 5G networks~\cite{10077211}, as it directly impairs latency-sensitive services such as video streaming and interactive communications applications.} We aim to accurately classify the underlying root causes based on KPIs patterns across different protocol layers.
High latency can result from various factors. {\color{R1color}The six root-cause categories considered in this work, as detailed in Table~\ref{tab:category_description}, represent the most commonly reported and operationally significant fault types in commercial mobile networks~\cite{klaine2017survey, ramirez2023explainable, yuan2020anomaly}.}

Let each sample $\mathbf{X} \in \mathbb{R}^{m \times l}$ denote the time-series record of a single degraded session, where $m$ is the total number of KPIs aggregated from the PHY, MAC, RLC, and PDCP layers, and $l$ is the temporal length of the monitoring window. Specifically, $\mathbf{X} = [\mathbf{x}_1, \mathbf{x}_2, \ldots, \mathbf{x}_m]^\top$, where each feature vector $\mathbf{x}_i \in \mathbb{R}^{l}$ represents the temporal evolution of the $i$-th KPI over $l$ consecutive intervals.


\begin{table*}
\centering
\caption{Target root-cause categories and their descriptions.}
\renewcommand{\arraystretch}{1.1}
\begin{tabular}{|c|c|c|}
\hline
\textbf{ID} & \textbf{Category} & \textbf{Description} \\
\hline
    1 & Uplink Interference & Severe interference detected on the uplink, potentially caused by external sources or frequency overlap. \\
    2 & Uplink Weak Coverage & Low uplink signal strength or quality, often due to poor coverage, penetration loss, or user location. \\
    3 & Downlink Interference & High interference on the downlink, leading to reduced decoding success at the user side. \\
    4 & Downlink Weak Coverage & Weak downlink signal reception caused by distance, obstruction, or poor cell coverage. \\
    5 & Traffic Channel Overload & Overload in the data channel, typically due to high traffic demand and resource exhaustion. \\
    6 & Control Channel Overload & Excessive load on control channels, affecting resource scheduling efficiency. \\
\hline
\end{tabular}
\label{tab:category_description}
\end{table*}

Two types of labeled samples are available for training, reflecting the dual nature of data-and-knowledge-driven supervision:

\FloatBarrier
\begin{itemize}
    \item \textbf{Rule-based labeled samples} $\mathcal{D}_r = \{(\mathbf{X}_r^{(i)}, y_r^{(i)})\}_{i=1}^{N_r}$: rule-based labels are automatically generated at scale by applying domain-specific thresholding rules to KPIs sequences. These labels reflect a data-driven paradigm, leveraging heuristics over large volumes of data but often suffering from noise and limited precision.
    \item \textbf{Expert-based labeled samples} $\mathcal{D}_e =\{(\mathbf{X}_e^{(i)}, y_e^{(i)})\}_{i=1}^{N_e}$: expert-based labels are carefully annotated by experienced communication domain experts based on deep domain knowledge. These labels, though limited in number, embody high-quality human expertise and serve as a knowledge-driven supervisory signal.
\end{itemize}



Overall, Figure~\ref{fig:bigpicture} summarizes the task setup. The goal is to learn a classification model $f: \mathbb{R}^{m \times l} \rightarrow {1, 2, \ldots, 6}$ that predicts the most likely root-cause label for a given KPIs time-series input $\mathbf{X}$, effectively leveraging both the abundant but noisy rule-based labels and the limited yet accurate expert annotations. This enables automatic and reliable analysis of root causes behind QoE degradations in mobile communication systems.


\section{Methodology}
% \section{Proposed method}
\label{sec:proposed method}

 

% First describe the problem addressed by the paper
\begin{figure}
    \centering
    \includegraphics[width=1\linewidth]{Figure/DK-Root.pdf}
    \caption{Overview of the DK-Root framework.}
    \label{fig: pipeline}
\end{figure}

This paper aims to address the task of classifying the root causes of communication-quality degradation when two sources of supervision are available: a small set of high-fidelity \emph{expert-based} labels and a large set of automatically generated \emph{rule-based} labels whose accuracy is less reliable. Figure~\ref{fig: pipeline} illustrates the proposed three-stage learning framework.

\noindent
\textbf{Stage I-Conditional Diffusion Training for KPIs Augmentation.}  
To leverage expert knowledge for distribution-aware augmentation, we train a label-conditioned diffusion model on the union of expert-labeled samples. Conditioning forces the model to learn class-specific data distributions, and because expert labels are noise-free, they act as an anchor that guides the diffusion process toward the true class manifold.

\noindent
\textbf{Stage II-Augmented KPIs-based Contrastive Representation Learning.}  
To construct a robust representation space under noisy supervision, the diffusion model is frozen and switched to inference mode. For every original sample, we generate two label-consistent views and apply a supervised contrastive loss that pulls embeddings of the same class together and pushes different classes apart, thereby yielding robust and class-discriminative representations without relying on cross-entropy over noisy rule labels.

\noindent
\textbf{Stage III-Model Fine-tuning for Root-Cause Classification.}  
To refine decision boundaries with high-fidelity labels, we attach a lightweight classification head to the encoder and fine-tune the model on the trusted expert-labeled subset, yielding an effective root-cause classifier.

The remainder of this section provides the implementation details of each stage.

\subsection{Stage I: Conditional Diffusion Training for KPIs Augmentation}
Diffusion models aim to model complex data distributions by gradually perturbing the input data with Gaussian noise and then learning to reverse this corruption process through a sequence of denoising steps. Among them, the Denoising Diffusion Probabilistic Model (DDPM)~\cite{ho2020denoising} has emerged as a foundational framework. 
It defines two stochastic processes: a forward diffusion process that adds Gaussian noise to the input data over multiple steps, progressively destroying information, and a reverse generative process that reconstructs the original data from noisy observations by learning the denoising transitions. The reverse process is parameterized by a neural network trained to predict the added noise at each timestep, which enables sample generation by iteratively denoising pure noise. This simple yet powerful training scheme allows DDPM to approximate highly complex, high-dimensional, and even multi-modal data distributions. 

While unconditional diffusion models are capable of capturing complex data distributions, they lack the ability to guide the generation process toward specific attributes or desired outcomes. To address this limitation, conditional diffusion models have been introduced~\cite{lee2024diffusionnag, dhariwal2021diffusion}. These models incorporate conditioning information (e.g., class labels, auxiliary features, or encoded context) into the reverse denoising steps, thereby enabling the generation of application-specific outputs tailored to task objectives. By injecting relevant external information, conditional diffusion improves controllability and semantic consistency, making it particularly suitable for tasks requiring structured and interpretable data synthesis~\cite{yang2024survey}.


\begin{figure}
    \centering
    \includegraphics[width=\linewidth]{Figure/diffusion-model.png}
    \caption{Conditional diffusion model for KPI-sequence augmentation.}
    \label{fig: flowchart_cond_diffusion}
\end{figure}

In our setting, each poor-QoE KPIs sample is denoted as $\mathbf{X} \in \mathbb{R}^{m \times l}$ (equivalently denoted as $\mathbf{X}^{(0)}$, representing the original sample before noise corruption) and the KPIs feature space is governed by a complex distribution $p(\mathbf{X})$, reflecting the intricate impact of diverse degradation causes across network protocol layers.  In this stage, we aim to learn this distribution using the conditional diffusion model illustrated in Figure~\ref{fig: flowchart_cond_diffusion}. Then, by sampling from the learned distribution, the model generates diverse yet semantically aligned augmented samples, which are subsequently used in the contrastive learning stage.


% Paragraph 2: describe model details
\paragraph{Forward Diffusion}

In the forward diffusion process, we gradually corrupt each clean KPIs time-series sample $\mathbf{X}^{(0)}$ by injecting Gaussian noise over $T$ diffusion steps. This process forms a Markov chain ${\mathbf{X}^{(t)}}_{t=1}^T$, where the distribution of the noisy sample at time step $t$ is given by:
\begin{equation}
q(\mathbf{X}^{(t)} \mid \mathbf{X}^{(0)}) = \mathcal{N}(\mathbf{X}^{(t)};\ \sqrt{\bar{\alpha}_t}, \mathbf{X}^{(0)},\ (1 - \bar{\alpha}_t), \mathbf{I}),
\end{equation}
where $\bar{\alpha}_t = \prod_{s=1}^t \alpha_s$ denotes the cumulative product of noise scaling factors. 

At each time step $t$, the clean input $\mathbf{X}^{(0)}$ is scaled by $\sqrt{\bar{\alpha}_t}$ and perturbed with zero-mean Gaussian noise of variance $(1 - \bar{\alpha}_t)$. This process gradually transforms the structured KPIs sequence into pure noise as $t \rightarrow T$, allowing the model to learn a denoising trajectory during reverse diffusion.

\paragraph{Class-Conditional Embedding and Fusion Module}

Each KPIs sample $\mathbf{X}^{(0)} \in \mathbb{R}^{m \times l}$ is associated with a discrete root-cause label $y$. 
To enable class-aware generative modeling, we incorporate degradation-type labels as conditioning information into the denoising network via a label embedding and early fusion strategy. 

Specifically, we adopt a soft embedding strategy\footnote{In practice, we also embed the diffusion timestep $t$ using an embedding layer and fuse it into the input in a similar manner (details omitted for brevity).}, where each label $y$ is mapped to a continuous vector through a trainable embedding module:
\begin{equation}
\mathbf{e}_y = \mathrm{Embed}_{\text{label}}(y) \in \mathbb{R}^d,
\end{equation}
followed by a linear projection to match the input channel dimension. Then, the projected label vector $\mathbf{e}_y'$ is broadcast along the temporal axis to form a condition tensor:
\begin{equation}
\mathbf{E}_y = \mathbf{e}_y' \otimes \mathbf{1}_l \in \mathbb{R}^{m \times l},
\end{equation}
where $\otimes$ denotes row-wise replication over $l$ time steps.

We concatenate the noisy input $\mathbf{X}^{(t)} \in \mathbb{R}^{m \times l}$ and the label embedding $\mathbf{E}_y$ along the channel dimension:
\begin{equation}
\tilde{\mathbf{X}}^{(t)} = \mathrm{Concat}\left(\mathbf{X}^{(t)},\ \mathbf{E}_y\right) \in \mathbb{R}^{2m \times l}.
\end{equation}
To effectively fuse the label condition with the input sequence, we apply a $1$-D convolutional layer:
\begin{equation}
\mathbf{X}_{\text{cond}}^{(t)} = \mathrm{Conv}(\tilde{\mathbf{X}}^{(t)}) \in \mathbb{R}^{m \times l},
\end{equation}
which projects the concatenated representation back to the original input dimension. 

The condition-aware representation $\mathbf{X}_{\text{cond}}^{(t)}$ is then fed into the U-Net\cite{ronneberger2015u} backbone to predict the noise $\hat{\boldsymbol{\epsilon}}_\theta(\mathbf{X}_{\text{cond}}^{(t)}, t, y)$ at timestep $t$.
During training, the diffusion model is optimized by minimizing the mean squared error (MSE) between the predicted noise and the true noise injected during the forward process:
\begin{equation}
\mathcal{L}_{\text{diff}} = \mathbb{E}_{\mathbf{X}^{(0)}, y, t, \boldsymbol{\epsilon}} \left[ \left\| \hat{\boldsymbol{\epsilon}}_\theta(\mathbf{X}_{\text{cond}}^{(t)}, t, y) - \boldsymbol{\epsilon} \right\|^2 \right],
\end{equation}

This early-stage fusion allows the network to extract class-specific features from the beginning, enabling label-conditioned signal encoding throughout the entire U-Net hierarchy while introducing minimal computational overhead compared with repeated conditioning in intermediate layers.

\paragraph{Reverse Process and Augmentation}
Once trained, the denoising network $\hat{\boldsymbol{\epsilon}}_\theta$ can be leveraged during inference to perform label-guided data augmentation. Unlike standard DDPMs that synthesize inputs from pure noise, our method begins with a real KPI sample $\mathbf{X}^{(0)}$ and generates a semantically aligned perturbation through a controlled diffusion-reversal process.

Specifically, we first apply a forward diffusion step to transform the original sample into its noisy version at timestep $t$, denoted as $\mathbf{X}^{(t)}$. The trained network then predicts the corresponding noise component $\hat{\boldsymbol{\epsilon}}_\theta(\mathbf{X}^{(t)}, t, y)$, conditioned on the degradation type $y$. Then, the augmented view is generated via a single-step reverse update:
\begin{equation}
\hat{\mathbf{X}}_{aug} = \frac{1}{\sqrt{\bar{\alpha}_t}} \left( \mathbf{X}^{(t)} - \sqrt{1 - \bar{\alpha}_t} \cdot \hat{\boldsymbol{\epsilon}}_\theta(\mathbf{X}^{(t)}, t, y) \right).
\end{equation}
where $\bar{\alpha}_t$ denotes the cumulative product of forward noise coefficients up to timestep $t$.

This data augmentation strategy yields semantically aligned views that preserve degradation type while introducing meaningful variations, enhancing representation learning and robustness in downstream tasks.

\subsection{Stage II: Augmented KPIs-based Contrastive Representation Learning}

% Figure: supervised contrastive learning workflow
\begin{figure}
    \centering
    \includegraphics[width=1\linewidth]{Figure/SupCon.png}
    \caption{Stage~II of DK-Root: contrastive representation learning with augmented KPI sequences.}
    \label{fig:enter-label}
\end{figure}

% Data augmentation strategy
The core strategy of contrastive learning is to minimize the distance between positive pairs constructed from augmented data. Undoubtedly, data augmentation plays a critical role in contrastive learning. Common augmentation techniques for time series can be broadly categorized into several classes.  
Transformation-based methods~\cite{um2017data, le2016data, goubeaud2021white} apply operations such as jittering, scaling, or magnitude warping directly to raw time-series signals to introduce variability.  Decomposition-based methods~\cite{wen2019robuststl, wen2020fast} decompose time series into subcomponents or segments (e.g., trend, seasonality), augment local timestamps or windows independently, and then reassemble them to construct novel sequences.
In contrast, generative-based methods~\cite{garcia2022improving, yang2023ts, chen2019emotionalgan} aim to capture the underlying data distribution using generative models such as Generative Adversarial Networks (GANs) or Variational Autoencoders (VAEs), from which new synthetic instances can be produced. However, when applied to communication quality degradation sequences, these models still face challenges in reproducing rare but critical anomalies that occur at specific timestamps, distorting critical indicators (e.g., RSRP, SINR) and thereby compromising the integrity of the root-cause analysis task. 


Traditional forecasting or generative models often make conservative predictions, constraining outputs within the historical value range, and thus struggle to generate low-probability abnormal patterns~\cite{lee2023maat}.  
To address this limitation, diffusion models are increasingly favored
% , as they offer a more expressive framework for generating diverse and semantically consistent abnormal patterns,
and are incorporated into the proposed framework. They learn to approximate the full data distribution by gradually denoising a Gaussian prior through a reverse diffusion process, enabling the modeling of not only the high-density regions but also the low-probability, anomalous patterns that are critical in communication quality degradation scenarios.  

However, in the task considered in this paper, rule-based labels inevitably contain noise, which can introduce bias into the generated samples if a standard diffusion model is applied directly. While diffusion models are predominantly employed for data generation, their probabilistic formulation also makes them well-suited for data augmentation, where synthetic samples are generated with additional information from the original data. This motivates us to design task-specific augmentation strategies that explicitly account for label imperfections and guide the generation process toward producing semantically consistent, label-aligned variations that enhance the downstream root-cause analysis objective.

\paragraph{Data Augmentation by Pretrained Diffusion Model}
With the conditional diffusion model trained in Stage~I, we perform semantic-preserving data augmentation on the weakly labeled dataset $\{(\mathbf{X}_r^{(i)})\}_{i=1}^{N_r}$, as illustrated in Figure~\ref{fig:enter-label}. 
The objective is to enrich the training set with diverse yet label-consistent variations, thereby improving the performance of representation learning.

Specifically, we introduce two levels of noise injection to create \emph{weak} and \emph{strong} augmentations. Let $T$ denote the total number of diffusion steps in the forward process. 
The augmentation strength is governed by the diffusion time step $t$, which controls the amount of noise added before the denoising process.
For each KPIs sample $\mathbf{X}_r$ in rule-based label dataset, we randomly sample two diffusion steps $t_{\mathrm{weak}}$ and $t_{\mathrm{strong}}$ from non-overlapping subranges of $[1, T]$:
\begin{equation}
t_{\mathrm{weak}} \sim \mathcal{U}(1, \alpha T), 
\quad 
t_{\mathrm{strong}} \sim \mathcal{U}(\beta T, T),
\label{for:alpha}
\end{equation}
where $0 < \alpha \leq \beta < 1$ are hyperparameters controlling the relative intensity of weak and strong augmentations. 
Typically, $\alpha$ is set to a small value to ensure minimal semantic distortion, while $\beta$ is set to a larger value to introduce more substantial perturbations.

After injecting noise according to the sampled $t$, the diffusion model performs reverse denoising to generate the augmented samples $\hat{\mathbf{X}}_{\mathrm{aug}}^{(t_{\mathrm{weak}})}$ and $\hat{\mathbf{X}}_{\mathrm{aug}}^{(t_{\mathrm{strong}})}$. 
These augmented views are expected to preserve the intrinsic semantics of the original KPIs sequence while introducing label-consistent variations that sharpen class boundaries across different root-cause categories.


\paragraph{Rule-based Label Aware Representation Learning}
We then train the encoder under a supervised contrastive learning paradigm to acquire preliminary discriminative ability.


Concretely, the two augmented views are passed through a shared encoder $f_\theta(\cdot)$ to obtain feature representations $\mathbf{z}_1$ and $\mathbf{z}_2$:
\[
\mathbf{z}^{(1)} = f_\theta\!\left(\hat{\mathbf{X}}_{\mathrm{aug}}^{(t_{\mathrm{strong}})}\right), \quad
\mathbf{z}^{(2)} = f_\theta\!\left(\hat{\mathbf{X}}_{\mathrm{aug}}^{(t_{\mathrm{weak}})}\right).
\]
Each representation is then flattened and $\ell_2$-normalized as
\[
\tilde{\mathbf{z}} = \frac{\mathrm{vec}(\mathbf{z})}{\lVert \mathrm{vec}(\mathbf{z}) \rVert_2},
\]
where $\mathrm{vec}(\cdot)$ denotes the vectorization operation.


Let $y_i$ denote the rule-based label assigned to the $i$-th KPIs sample. 
In our supervised contrastive framework, a positive pair is defined as any two samples $(i,j)$ satisfying $y_i = y_j$ and $i \neq j$. 
We use the indicator function $\mathbb{I}[\cdot]$ to represent the positive-pair relation:
\[
M_{ij} = \mathbb{I}[y_i = y_j], \qquad M \in \{0,1\}^{N_r \times N_r}.
\]

Each sample is augmented twice to produce two distinct views, which are encoded and $\ell_2$-normalized to form the contrastive feature set:
\[
\mathcal{C} = \left[ \tilde{\mathbf{z}}^{(1)}_1, \dots, \tilde{\mathbf{z}}^{(1)}_{N_r}, 
\tilde{\mathbf{z}}^{(2)}_1, \dots, \tilde{\mathbf{z}}^{(2)}_{N_r} \right] 
\in \mathbb{R}^{2N_r \times D}.
\]
where $D$ is the dimensionality of the $\ell_2$-normalized embeddings.

The similarity between two feature vectors $\mathbf{a}$ and $\mathbf{c}$ is measured as
\[
s(\mathbf{a},\mathbf{c}) = \frac{\mathbf{a}^\top \mathbf{c}}{\tau},
\]
where $\tau>0$ is the temperature parameter.

Let $s_{ik} = s(\tilde{\mathbf{z}}_i, \tilde{\mathbf{z}}_k)$ 
denote the pairwise similarity between anchor $i$ and contrast feature $k$.
Then, the log-probability is given by the log-softmax over all non-anchor features:
\[
\log p_{ik} = \log \frac{\exp(s_{ik})}{\sum_{k' \neq i} \exp(s_{ik'})}.
\]
Let $P(i) = \{ k \,:\, M_{ik} = 1 \}$ be the positive set for anchor $i$. 
The supervised contrastive loss is defined as
\begin{equation}
\mathcal{L}_{con} = - \frac{1}{2N_r} \sum_{i=1}^{2N_r} 
\frac{1}{|P(i)|} \sum_{k \in P(i)} \log p_{ik}.
\end{equation}
This objective encourages representations of KPI samples sharing the same rule-based label to be pulled closer, while pushing apart those with different root-cause classes, thereby promoting class-discriminative feature learning. The resulting encoder acquires preliminary discriminative ability, which can be further calibrated in a subsequent fine-tuning stage using scarce but high-quality expert annotations.



\subsection{Stage III: Model Fine-tuning for Root-Cause Classification}

The third stage is dedicated to the final degradation classification task. 
While Stage~II produces class-discriminative representations under noisy rule-based supervision, the purpose of this stage is to adapt these representations to high-precision decision boundaries using only clean expert annotations. 
This separation ensures that the learned encoder is first guided by abundant but potentially noisy data, and is then refined using scarce yet reliable labels to maximize classification accuracy.

We use the expert-labeled dataset 
$\mathcal{D}_e = \{(\mathbf{X}_e^{(i)}, y_e^{(i)})\}_{i=1}^{N_e}$ 
for supervised training. 
A lightweight classification head is attached to the encoder obtained from Stage~II, and the entire model is fine-tuned end-to-end to predict degradation categories. 
The classification head outputs class probabilities via a softmax layer:
\begin{equation}
p_c^{(i)} = \frac{\exp(h_c^{(i)})}{\sum_{c'=1}^6 \exp(h_{c'}^{(i)})},
\end{equation}
where $h_c^{(i)}$ is the logit score for class $c$.

We adopt the standard multi-class cross-entropy loss:
\begin{equation}
\mathcal{L}_{sup} = -\frac{1}{N_e} \sum_{i=1}^{N_e} \sum_{c=1}^6 
\mathbb{I}[y_e^{(i)} = c] \, \log p_c^{(i)},
\end{equation}


\section{Experiments}
\label{sec:experiments}

\subsection{Dataset and Setup}

The dataset used in this study was collected by Huawei from real-world mobile-network sensors at 5-second intervals, resulting in multivariate time-series samples with a fixed sequence length of 40, i.e., 200 seconds per sample. The dataset contains {\color{R3color}3,000} samples with rule-based labels, which are automatically generated by applying structured rules to KPIs, and {\color{R3color}375} samples with manually annotated expert labels.
A detailed list of KPIs, the per-class distribution of the expert-labeled dataset, and a representative example of the rule-based labeling strategy are provided in Appendix~\ref{appendix}.
Classification accuracy is used as the primary evaluation metric. {\color{AEcolor}We additionally report macro-averaged F1-score and confusion matrices to provide a fine-grained view of performance across all six root-cause categories. Formal definitions of all metrics are provided in Appendix~\ref{appendix:metrics}.}
{\color{R3color}
% Table~\ref{tab:implementation} summarizes the key architectural and training hyperparameters for each stage. 
All experiments are conducted on a single NVIDIA RTX 3090 GPU using PyTorch.
}


We conduct six sets of experiments to comprehensively evaluate the proposed framework:
\begin{itemize}
    \item To verify the overall effectiveness of \textit{DK-Root} and provide a systematic understanding of root-cause analysis in realistic dual-label settings, we conduct baseline comparisons against a diverse set of classical, semi-supervised, and {\color{R3color}weakly-supervised} methods (Section~\ref{sec:rca_effectiveness}).
    \item {\color{R3color}To investigate the interplay between data regime and model behavior, and to reveal how each label source contributes to diagnostic capability, we conduct sensitivity analyses on expert-label and rule-label quantities, yielding deeper insights into QoE anomaly characteristics (Section~\ref{sec:understanding}).}
    \item {\color{R1color}To assess the framework's resilience against imperfect supervision, we conduct controlled noise injection experiments by artificially corrupting rule-based labels at varying rates, examining the degradation boundary of the contrastive pretraining stage (Section~\ref{sec:noise_robustness}).}
    \item {\color{R2color}To validate the effectiveness of the augmentation strategy, we analyze sensitivity to key hyperparameters and evaluate the semantic consistency of augmented samples, comparing the proposed method with several traditional and generative alternatives (Section~\ref{sec:diffusion_effectiveness}).}
    \item {\color{AEcolor}To reveal the structural benefit of the multi-stage design on per-class discriminability, we conduct ablation studies including latent-space visualization, convergence comparison, and confusion-matrix analysis across different training paradigms (Section~\ref{sec:pretrain_finetune}).}
    \item {\color{R2color} We finally report the offline training cost and online inference latency of the proposed framework (Section~\ref{sec:deployment}).}
\end{itemize}

\subsection{Overall Identification Performance}
\label{sec:rca_effectiveness}


\noindent\textbf{Baselines.}  
We compare \textit{DK-Root} against three groups of methods to observe whether \textit{DK-Root} improves root-cause identification under realistic dual-label supervision: classical ML methods (Sliding-KNN, SVM, XGBoost, TSF, BOP), semi-supervised methods (\textit{TS-TCC}~\cite{eldele2021time}, \textit{CA-TCC}~\cite{eldele2023self}), and {\color{R3color}weakly-supervised methods (\textit{Co-teaching}~\cite{han2018coteaching}, \textit{DivideMix}~\cite{li2020dividemix}, \textit{Scale-teaching}~\cite{liu2023scale})}. Detailed descriptions of all baselines are provided in Appendix~\ref{appendix:baselines}.



\begin{table}
\centering
\caption{{\color{AEcolor}Comparison of root-cause identification performance across DK-Root and baseline methods.}}
\resizebox{0.48\textwidth}{!}{
\begin{tabular}{ll cc}
\toprule
\textbf{Category} & \textbf{Method} & \textbf{Accuracy} & {\color{AEcolor}\textbf{Macro-F1}} \\
\midrule
\multirow{5}{*}{{ML}}
& SVM             & 0.4737$\pm$0.0220 & 0.3896$\pm$0.0156 \\
& XGBoost         & 0.4500$\pm$0.0337 & 0.4082$\pm$0.0358 \\
& Sliding-KNN     & 0.2855$\pm$0.0198 & 0.2145$\pm$0.0198 \\
& BOP             & 0.2579$\pm$0.0229 & 0.1682$\pm$0.0285 \\
& TSF             & 0.3592$\pm$0.0219 & 0.2805$\pm$0.0200 \\
\midrule

\multirow{2}{*}{Semi-SL}
& TS-TCC          & 0.5792$\pm$0.0387 & 0.5187$\pm$0.0369\\
& CA-TCC          & 0.5250$\pm$0.0081 & 0.4033$\pm$0.0278 \\
\midrule
{\color{R3color}
\multirow{3}{*}{{Weakly-SL}}}
& {\color{R3color}Co-teaching}     & {\color{R3color}0.6377$\pm$0.0383} & {\color{R3color}0.5990$\pm$0.0348} \\
& {\color{R3color}DivideMix}       & {\color{R3color}0.6518$\pm$0.0374} & {\color{R3color}0.6136$\pm$0.0382} \\
& {\color{R3color}Scale-teaching}   & {\color{R3color}0.6566$\pm$0.0459} & {\color{R3color}0.6363$\pm$0.0488} \\
\midrule
& \textbf{DK-Root} & \textbf{0.6907$\pm$0.0310} & \textbf{0.6867$\pm$0.0344} \\
\bottomrule
\end{tabular}
\label{tab:baseline_comparison}
}
\end{table}


\noindent\textbf{Overall Comparison.} Table~\ref{tab:baseline_comparison} shows that \textit{DK-Root} achieves the best accuracy and macro-F1 among all compared methods, while also maintaining the smallest gap between the two metrics. This pattern indicates not only stronger average performance but also more balanced behavior across categories. Among the classical ML baselines, SVM and XGBoost provide moderate results, whereas representation-learning-based methods consistently outperform them. {\color{R3color}Compared with semi-supervised methods that rely only on the limited expert-labeled set, all weakly supervised baselines and \textit{DK-Root} perform better, illustrating that the additional supervision provided by rule-based labels remains useful. Moreover, \textit{DK-Root} performs better possibly because rule-based labels are mainly exploited to improve representations during contrastive pretraining, while the final classifier is refined using clean expert labels. By contrast, weakly supervised baselines rely more directly on noisy rule-based supervision during classifier learning, even when sample selection, semi-supervised training, or label correction mechanisms are used to mitigate label noise.}
% By contrast, weakly supervised baselines still learn decision boundaries more directly from filtered noisy labels through cross-entropy training.} 


{\color{AEcolor}
{\color{R2color}\noindent\textbf{Advantages in Mitigating Root-Cause Disambiguation.}
Figure~\ref{fig:combined_performance} reveals that QoE root-cause disambiguation remains intrinsically difficult in real networks. The dominant confusion appears between \textit{UL Interference} and \textit{UL Weak Coverage} because both causes degrade the same uplink signal-quality indicators (RSRP and SINR) even though they originate from different physical mechanisms, namely co-channel interference and geometric path loss. By comparison, overload-related categories are more separable for all methods because PRB utilization and CCE scheduling expose clearer signatures of resource exhaustion. 
We observe that \textit{DK-Root} achieves more balanced performance across categories, especially for \textit{DL Weak Coverage} and \textit{Traffic Channel Overload}, where competing weakly-supervised methods degrade substantially, suggesting that the proposed representation learning stage captures distinctions that are easy to blur under direct noisy-label supervision.}
}

\begin{figure}
    \centering
    \begin{subfigure}{\linewidth}
        \centering
        \includegraphics[width=\linewidth]{Figure/confusion_matrix_2x2.pdf}
        \caption{Confusion matrices of DK-Root and three weakly supervised baselines.}
        \label{fig:confusion_matrix}
    \end{subfigure}
    
    \begin{subfigure}{\linewidth}
        \centering
        \includegraphics[width=\linewidth]{Figure/per_class_acc_comparison.pdf}
        \caption{Per-class accuracy comparison across DK-Root and three weakly supervised baselines.}
        \label{fig:per_class_acc}
    \end{subfigure}
    \caption{Class-level performance comparison between DK-Root and three weakly supervised baselines.}
    \label{fig:combined_performance}
\end{figure}


\subsection{Impact of Label Quantity}
\label{sec:understanding}

\begin{figure}
    \centering
    \begin{subfigure}[b]{1\linewidth}
        \centering
        \includegraphics[width=\linewidth]{Figure/ne_sensitivity.pdf}
        \caption{Identification performance versus the number of expert-labeled samples $N_e$.}
        \label{fig:ne_sensitivity}
    \end{subfigure}
    
    \begin{subfigure}[b]{1\linewidth}
        \centering
        \includegraphics[width=\linewidth]{Figure/nr_sensitivity.pdf}
        \caption{Identification performance versus the number of rule-labeled samples $N_r$.}
        \label{fig:nr_sensitivity}
    \end{subfigure}
    \caption{Sensitivity of identification performance to the number of labeled samples.}
    \label{fig:sensitivity_combined}
\end{figure}

{\color{R3color}
This subsection answers how the amounts of expert-labeled and rule-labeled data affect identification performance. To this end, we conduct two sensitivity analyses: one varies $N_e$ while keeping the rule-labeled set fixed, and the other varies $N_r$ while keeping $N_e$ fixed.


\noindent\textbf{Benefit of Incorporating Rule-Based Labels.}
Figure~\ref{fig:sensitivity_combined}(a) shows that expert labels remain the primary driver of final classification performance, but \textit{DK-Root} consistently outperforms the purely supervised baseline across all values of $N_e$. The gain is most pronounced in the low-label regime, where representation learning from noisy but abundant rule-labeled data is particularly valuable. As $N_e$ increases, the gap gradually narrows, indicating that sufficient expert supervision can eventually compensate for the absence of pretraining. More importantly, \textit{DK-Root} reaches a performance level close to the supervised baseline with substantially fewer expert labels. For example, the curve suggests comparable performance around $N_e=50$ versus $100$ and around $N_e=225$ versus $340$. This reduction is practically meaningful because expert annotation is labor-intensive, and each expert-labeled sample typically requires nontrivial diagnosis time.

\noindent\textbf{Sensitivity to Rule-Label Quantity.}
Figure~\ref{fig:sensitivity_combined}(b) shows that \textit{DK-Root} is robust to the quantity of rule-based labels. Performance improves rapidly once the model has access to the first few hundred rule-labeled samples, then the curve gradually plateaus and mild oscillations appear when more rule-labeled samples are added. The curve illustrates that a larger weakly labeled pool also introduces more noisy and potentially conflicting supervision signals. Nevertheless, the performance remains consistently above the purely supervised baseline over a wide range of $N_r$, showing that the benefit of rule-based labels is stable rather than fragile. Taken together, these results confirm that expert labels determine the final decision boundary, whereas rule-labeled data mainly improves the representation quality that supports that boundary.

% These results collectively provide a \textit{data-centric understanding} of the QoE anomaly classification problem. the six root-cause categories exhibit separable but nuanced KPI signatures across protocol layers, and the primary bottleneck in diagnosis is not the quantity of either label source alone, but rather the synergistic use of noisy-but-abundant rule labels for representation learning and scarce-but-precise expert labels for decision boundary refinement.

}

\subsection{Robustness to Noisy Rule Labels}
\label{sec:noise_robustness}
{\color{R1color}
To evaluate whether heavily corrupted rule-based labels can negate the benefit of Stage~II pretraining, we artificially inject additional noise into the rule-based labels before Stage~II training, pushing the total noise rate well beyond the empirically observed 27.4\% disagreement rate between rule-based and expert labels.

\begin{figure}
    \centering
    \includegraphics[width=1\linewidth]{Figure/noise_robustness.pdf}
    \caption{{\color{R1color}Robustness of DK-Root to injected rule-label noise. The dashed horizontal line indicates the supervised baseline.}}
    \label{fig:noise_robustness}
\end{figure}

Figure~\ref{fig:noise_robustness} reveals a key robustness property of \textit{DK-Root}. In Stage~II, the contrastive learning covers two types of positive pairs: \emph{cross-sample pairs}, formed between different samples that share the same rule-based label, and \emph{intra-instance pairs}, formed between the two augmented views ($\hat{\mathbf{X}}^{(t_\mathrm{weak})}$ and $\hat{\mathbf{X}}^{(t_\mathrm{strong})}$) of the same sample. As rule-label noise increases, cross-sample pairs become less reliable because they increasingly connect semantically unrelated samples. By contrast, intra-instance pairs remain valid because the Stage~I diffusion model is trained only on clean expert labels, so the two augmented views of a sample remain semantically consistent with the original input. In the extreme case of heavily corrupted rule labels, CL therefore degrades toward instance-level pretraining. Stage~III fine-tuning with clean expert labels then restores class boundaries on top of this coherent representation, which explains why the proposed framework remains competitive even under severe rule-label corruption.

These results indicate that different noise levels in the introduced rule labels affect the amount of performance gain, but they do not destroy the core representation-learning benefit brought by the pretraining stage. This observation is also consistent with the trends reported in Section~\ref{sec:understanding}.
}


\begin{table*}
\centering
\caption{{\color{R2color}Comparison of semantic consistency and class structure under different data augmentation methods.}}
\label{tab:augmentation_metrics}
\resizebox{\textwidth}{!}{
\begin{tabular}{l|c|ccc|c|cccccc}
\hline
\textbf{Method} 
& \textbf{Avg. Inter-class Dist. $\uparrow$} 
& \textbf{Silhouette $\uparrow$} 
& \textbf{CH $\uparrow$} 
& \textbf{DB $\downarrow$} 
& \textbf{MI $\uparrow$} 
& \textbf{Class1 $\downarrow$} 
& \textbf{Class2 $\downarrow$} 
& \textbf{Class3 $\downarrow$} 
& \textbf{Class4 $\downarrow$} 
& \textbf{Class5 $\downarrow$} 
& \textbf{Class6 $\downarrow$} \\
\hline
Original     & 5.1344 & 0.0255 & 4.7111 & 4.1954 & 0.5191 & 49.0086 & 47.1370 & 47.4579 & 26.3059 & 40.7622 & 43.4905 \\
DK-Root(strong)  & 4.4538 & 0.0171 & 4.8761 & 3.7663 & 0.5416 & 35.3691 & 26.8484 & 38.8719 & 25.0625 & 31.3818 & 34.0451 \\
DK-Root(weak)   & 4.9249 & 0.0244 & 4.7072 & 3.9884 & 0.4799 & 46.7216 & 47.3592 & 46.4654 & 28.6535 & 37.9558 & 42.0162 \\
{\color{R2color}TimeGAN(strong)}  & {\color{R2color}2.5980} & {\color{R2color}$-$0.1105} & {\color{R2color}1.0658} & {\color{R2color}7.4426} & {\color{R2color}0.1579} & {\color{R2color}35.6720} & {\color{R2color}27.5796} & {\color{R2color}38.0567} & {\color{R2color}30.4851} & {\color{R2color}19.7989} & {\color{R2color}38.3882} \\
{\color{R2color}TimeGAN(weak)}   & {\color{R2color}2.3559} & {\color{R2color}$-$0.1251} & {\color{R2color}0.9446} & {\color{R2color}6.9523} & {\color{R2color}0.1897} & {\color{R2color}31.0021} & {\color{R2color}18.2805} & {\color{R2color}30.7098} & {\color{R2color}22.9885} & {\color{R2color}29.3276} & {\color{R2color}31.0759} \\
Noise injection  & 27.1371 & $-$0.1075 & 0.8072 & 6.4504 & 0.2472 & 4460.90 & 3949.16 & 5058.29 & 2610.70 & 4818.82 & 6781.01 \\
Scaling      & 5.1627 & 0.0250 & 4.6104 & 4.2050 & 0.5054 & 49.7388 & 48.0329 & 48.4099 & 28.2337 & 42.0635 & 44.6788 \\
\hline
\end{tabular}
}
\end{table*}

\subsection{Effectiveness of Diffusion Augmentation}
\label{sec:diffusion_effectiveness}

In this subsection, we evaluate the effectiveness of the proposed diffusion-based augmentation strategy in Stage~II. Specifically, we examine whether the generated views preserve root-cause semantics while improving intra-class compactness and class separability. We compare \textit{DK-Root} with conventional augmentation methods and {\color{R2color}\textit{TimeGAN}~\cite{yoon2019timegan}}.\footnote{Detailed descriptions of all augmentation baselines and the evaluation metrics, namely Average Inter-class Distance, Silhouette Score, CH, DB, MI, and Intra-class Variance, are provided in Appendix~\ref{appendix:diffusion_baselines}.}

% We first quantify the perturbation strength introduced at different diffusion time steps, and then evaluate the augmented data distributions under multiple clustering and information-theoretic metrics. {\color{R2color}To further justify the choice of diffusion models over alternative generative approaches, we include a head-to-head comparison with TimeGAN~\cite{yoon2019timegan}




\begin{figure*}
    \centering
    \includegraphics[width=0.75\linewidth]{Figure/ws_diff_l2_dual.pdf}
    \caption{$L_2$ distance between original samples and augmented samples versus the diffusion step ratio.}
    \label{fig:hyperpara_l2}
\end{figure*}

{\color{R1color}
\noindent\textbf{Sensitivity to Diffusion Step Ranges.}
The hyperparameters $\alpha$ and $\beta$ in Eq.~(\ref{for:alpha}) control the intensity of weak and strong augmentation, respectively. To justify their selection, we sweep them and report the $L_2$ distances between the noised input and the original sample (Noise~$L_2$) and between the denoised output and the original sample (Denoise~$L_2$) for each setting.
As shown in Figure~\ref{fig:hyperpara_l2}, the two augmentation regimes exhibit clearly different perturbation strengths while both remaining recoverable after denoising. For weak augmentation, both Noise~$L_2$ and Denoise~$L_2$ remain relatively small across the full sweep, confirming that small $\alpha$ values produce conservative perturbations. For strong augmentation, the Noise~$L_2$ is roughly an order of magnitude larger, whereas the Denoise~$L_2$ grows much more slowly. 

{\color{R2color}
\noindent\textbf{Semantic Consistency Analysis.}
Table~\ref{tab:augmentation_metrics} evaluates augmentation quality from three angles: cluster geometry, semantic fidelity, and within-class diversity. \textit{DK-Root(strong)} achieves the best CH score, the lowest DB value, and consistently lower within-class variances than the original data, while also slightly improving MI ($0.5416$ vs.~$0.5191$). Although its average inter-class distance and Silhouette score are slightly below those of the original set, the overall pattern still indicates a better compactness--separability trade-off rather than semantic drift. By comparison, \textit{DK-Root(weak)} stays close to the original distribution on most metrics, which is desirable for a stable anchor view in contrastive learning.
Scaling changes the data distribution only marginally, so it provides limited additional training signal. Noise injection produces the largest average inter-class distance, but this value is misleading because it is accompanied by a negative Silhouette score, much lower CH, higher DB and sharply reduced MI. In other words, the perturbation mainly scatters samples instead of preserving meaningful class structure.
Both \textit{TimeGAN(strong)} and \textit{TimeGAN(weak)} yield negative Silhouette scores, low CH, high DB, and MI values, revealing severe semantic drift and boundary collapse. Based on these results, their relatively small within-class variances in some categories mean that GAN-based generation becomes over-smoothed or partially mode-collapsed, concentrating samples around a few generic prototypes while weakening label-specific distinctions. 
% For contrastive learning, this combination is doubly harmful: intra-instance positive pairs lose meaningful variation between weak and strong views, and cross-sample positive pairs drawn from the same nominal label may actually span multiple true semantic classes, injecting corrupted supervision into the contrastive objective.

These observations are consistent with the downstream classification results. When \textit{TimeGAN} augmentation is plugged into the same \textit{DK-Root} pipeline (Stage~II $\rightarrow$ Stage~III), it yields an accuracy of $0.6683$ and a macro-F1 of $0.6340$, both lower than those of diffusion-augmented \textit{DK-Root} in Table~\ref{tab:baseline_comparison}. Its performance is also close to the purely supervised setting in Section~\ref{sec:pretrain_finetune}, which suggests that semantically drifted positive pairs provide limited representation-learning benefit.

% Overall, the evidence shows that effective augmentation for this task must balance three properties at once: semantic consistency, sufficient diversity, and controllable perturbation strength. The proposed diffusion-based design is more aligned with this objective than either generic perturbation or GAN-based generation. 
}



% % Loss curve of Stage I in the proposed method
% \begin{figure}
%     \centering
%     \includegraphics[width=0.85\linewidth]{Figure/self_supervised_loss_curve_57.png}
%     \caption{Training loss curve of contrastive learning.}
%     \label{fig:loss_curve}
% \end{figure}


\begin{figure*}
    \centering
    \includegraphics[width=0.75\linewidth]{Figure/tsne.pdf}
    \caption{{\color{R1color}t-SNE visualizations of latent representations after Stage~II pretraining and Stage~III fine-tuning.}}
    \label{fig:tsne_comparison}
\end{figure*}

\begin{figure}[!t]
    \centering
    \includegraphics[width=0.75\linewidth]{Figure/ablation_curve.pdf}
    \caption{Learning curves of the three training paradigms.}
    \label{fig:acc_curve}
\end{figure}

\begin{figure*}[!t]
    \centering
    \includegraphics[width=0.95\linewidth]{Figure/confusion_matrix_1x3.pdf}
    \caption{Confusion matrices of the three training paradigms.}
    \label{fig:confusion_matrix_1x3}
\end{figure*}

% \newlength{\ablationheight}
% \settoheight{\ablationheight}{\includegraphics[width=0.74\linewidth]{Figure/confusion_matrix_1x3.pdf}}


\subsection{Contribution of Pretraining and Fine-tuning}
\label{sec:pretrain_finetune}

\noindent\textbf{Effect of Pretraining and Fine-tuning.}  
Figure~\ref{fig:tsne_comparison} visualizes the latent feature space using t-SNE. After Stage~II contrastive pretraining, the samples form rough class-level clusters, although noticeable inter-class overlap remains, especially between \textit{UL Weak Coverage} and \textit{UL Interference}. After Stage~III fine-tuning, the latent space becomes more structured and discriminative, with clearer inter-class margins. This progression shows that contrastive pretraining provides a useful initialization, whereas expert-guided fine-tuning sharpens the decision boundaries needed for final classification.


\noindent\textbf{Comparison of Different Training Strategies.}  
We also evaluate \textit{DK-Root} against \textit{CNN-exp} (the same as the \textit{supervised} baseline in Section~\ref{sec:noise_robustness}), a purely supervised CNN trained exclusively on expert-labeled data, and \textit{CNN-full}, a purely supervised CNN trained on the union of the expert-labeled and rule-labeled datasets\footnote{All CNN-based models share the same network architecture, identical to the encoder used in our proposed method.}. 
Figure~\ref{fig:acc_curve} shows that \textit{DK-Root} stabilizes at a clearly higher plateau. Notably, \textit{CNN-full} performs worse than the expert-only \textit{CNN-exp}, showing that naively mixing noisy rule-based labels into supervised training corrupts the decision boundary and degrades generalization. In contrast, using rule-based labels through contrastive learning produces more robust and noise-tolerant representations. This comparison shows that the benefit of the proposed pipeline comes from using noisy rule labels in a representation-learning stage that is less sensitive to label errors before the final classifier is refined with clean expert supervision.

{\color{AEcolor}

\noindent\textbf{Per-Class Performance across Training Paradigms.}
The confusion matrices in Figure~\ref{fig:confusion_matrix_1x3} show a clearer class-level picture. \textit{DK-Root} shows better balanced performance. Meanwhile, the difference appears on hard categories. For \textit{Traffic Channel Overload}, \textit{DK-Root} reaches $67\%$ recall, compared with only $42\%$ for \textit{CNN-exp} and $58\%$ for \textit{CNN-full}. On \textit{UL Interference}, \textit{CNN-exp} misclassifies $21\%$ of samples as \textit{Control Channel Overload}, indicating that the proposed pipeline better separates interference-related and overload-related patterns. These results suggest that contrastive pretraining on rule-labeled data is particularly useful for categories whose KPI signatures partially overlap, because it improves representation structure before the final expert-guided fine-tuning stage.
% In summary, the ablation results show that the proposed pipeline improves not only average accuracy but also class-level reliability on the categories that matter most.

}

\subsection{Deployment Cost and Practicality}
\label{sec:deployment}
{\color{R2color}To assess practical deployability, we report the wall-clock time of each stage in Table~\ref{tab:runtime}. Stages~I, II, and III are executed offline before deployment, whereas online inference uses only the CNN encoder and the lightweight linear classification head.}
Conditional diffusion training constitutes the main offline cost. At deployment time, however, the final model contains 26,770 trainable parameters and achieves sub-millisecond inference latency on GPU. Such latency is compatible with real-time network management systems that typically operate on minute-level KPI aggregation windows.

{\color{R2color}
\begin{table}[!b]
\centering
\caption{{\color{R2color}Offline training cost and online inference latency of DK-Root.}}
\label{tab:runtime}
\begin{tabular}{lcc}
\toprule
\textbf{Stage} & \textbf{Wall-Clock Time} & \textbf{Mode} \\
\midrule
Stage I: Diffusion training        & $411$~s         & Offline  \\
Stage I: Sample generation         & $0.09$~s       & Offline  \\
Stage II: Contrastive pretraining  & $105$~s          & Offline  \\
Stage III: Expert fine-tuning      & $33$~s          & Offline  \\
\midrule
Online inference/per sample &       $0.42$~ms      & Online \\
\bottomrule
\end{tabular}
\end{table}
}


\section{Conclusion}
\label{conclusion}

This paper presents \textit{DK-Root}, a three-stage semi-supervised framework for root-cause identification of communication quality degradation. By combining conditional diffusion augmentation, supervised contrastive pretraining, and fine-tuning on scarce expert labels, \textit{DK-Root} effectively unifies noisy rule-based labels with limited expert annotations. Experiments on a real-world operator-grade KPI dataset confirm that the framework achieves competitive accuracy and strong robustness across varying label quantities and noise levels.
{\color{R1color}Yet, current formulation assumes a single dominant root cause per case, whereas in practice multiple degradation mechanisms may co-occur. The same framework is also applicable to other QoE degradation types, such as low throughput and high packet loss, provided that the corresponding cross-layer KPIs can be collected and annotated. However, sufficiently large real-world datasets for these scenarios remain limited at present. Extending the framework to multi-root-cause analysis, broader QoE degradation types, and cross-network generalization is therefore an important direction for future work.}

{
\FloatBarrier
\appendices

\section{Dataset, Setup, KPI Definitions, and Rule-Based Labeling}
\label{appendix}

\subsection{Dataset Summary}

The experiments use an operator-grade KPI dataset with two supervision sources: a large rule-labeled set and a smaller expert-labeled set. Within the expert-labeled subset, the per-class distribution is UL Interference~(58), UL Weak Coverage~(84), DL Interference~(81), DL Weak Coverage~(21), Traffic Channel Overload~(28), and Control Channel Overload~(103). 

\subsection{Setup Details}

Table~\ref{tab:implementation} summarizes the main implementation and training settings used in the three stages of \textit{DK-Root}.

\begin{table}[!t]
\centering
\caption{{\color{R3color}Implementation details of each stage in DK-Root.}}
\label{tab:implementation}
\scriptsize
\begin{tabular}{l|ccc}
\toprule
    \textbf{Parameter} & \textbf{Stage I} & \textbf{Stage II} & \textbf{Stage III} \\
\midrule
Backbone      & U-Net   & 1D-CNN  & 1D-CNN \\
Optimizer     & Adam    & Adam    & Adam \\
Learning rate & 1e-3    & 3e-4    & 3e-4 \\
Batch size    & 8       & 16      & 16 \\
Epochs        & 1000    & 60      & 150 \\
\bottomrule
\end{tabular}
\end{table}

\subsection{KPI Definitions}

Table~\ref{tab:selected_features} summarizes the KPIs used as model-input features and/or rule-labeling signals together with their protocol-layer mapping.

\begin{table*}[!t]
\scriptsize
\caption{Selected KPIs and their protocol-layer mapping.}
\label{tab:selected_features}
\begin{tabular}{@{}>{\raggedright\arraybackslash}p{0.24\textwidth}|>{\raggedright\arraybackslash}p{0.65\textwidth}|>{\centering\arraybackslash}p{0.06\textwidth}@{}}
\hline
\textbf{Feature Name} & \textbf{Description} & \textbf{Layer} \\
\hline
{UL\_DMRS\_SINR} & SINR of the Uplink Demodulation Reference Signal (DMRS), indicating uplink signal quality & PHY \\
{UL\_WB\_PRE\_SINR} & Wideband pre-SINR measured at the uplink side before decoding, reflecting overall uplink signal quality & PHY \\
{UL\_SRS\_RSRP} & Received Signal Received Power (RSRP) of Uplink Sounding Reference Signal (SRS), used for uplink channel quality estimation & PHY \\
{UL\_DMRS\_RSRP\_AVG} & Average RSRP of Uplink DMRS, indicating average received power of uplink reference signal & PHY \\
{SS\_SINR} & SINR of the Synchronization Signal (SS) of the serving cell, representing downlink synchronization quality & PHY \\
{RSRP\_DIFF\_NEIGH\_CNT} & Number of same-frequency neighboring cells whose RSRP differs from the serving cell by more than a specified threshold (Thd1) & PHY \\
{TOP1\_NEIGH\_SSB\_RSRP} & RSRP of the strongest same-frequency neighboring SSB & PHY \\
{SERV\_SSB\_RSRP} & RSRP of the Serving Cell Synchronization Signal Block (SSB) & PHY \\
{SERV\_SSB\_RSRQ} & RSRQ (Reference Signal Received Quality) of the Serving Cell SSB & PHY \\
{CQI\_AVG\_CW0} & Average Channel Quality Indicator (CQI) for codeword 0, used for downlink modulation and coding scheme selection & PHY \\
{UL\_RBLER} & Uplink block error rate & PHY \\
{DL\_RBLER} & Downlink block error rate & PHY \\
{UL\_DMRS\_RSRP\_MIN} & Minimum RSRP of uplink DMRS signal & PHY \\
{DL\_PRB\_UTIL} & Downlink Physical Resource Block (PRB) utilization ratio of the cell & MAC \\
{UL\_DTX\_Ratio} & Ratio of uplink discontinuous transmission (DTX) & MAC \\
{UL\_MAC\_HARQ\_RETX\_MAX} & Maximum number of uplink HARQ retransmissions & MAC \\
{DL\_MAC\_HARQ\_RETX\_MAX} & Maximum number of downlink HARQ retransmissions & MAC \\
{UL\_PRB\_UTIL} & Uplink Physical Resource Block (PRB) utilization ratio of the cell & MAC \\
{DL\_CCE\_FAIL\_RATE} & Failure rate of downlink Control Channel Element (CCE) allocation & MAC \\
{UL\_CCE\_FAIL\_RATE} & Failure rate of uplink CCE allocation & MAC \\
{CCE\_UTIL\_RATE} & Overall utilization rate of CCEs in the cell & MAC \\
{DL\_CCE\_USAGE\_RATIO} & Proportion of CCEs used for downlink transmission & MAC \\
{COMMON\_CCE\_USAGE} & Proportion of common CCEs shared by control and data channels & MAC \\
{UL\_CCE\_USAGE\_RATIO} & Proportion of CCEs used for uplink transmission & MAC \\
{UL\_SCHED\_FAIL\_RATE\_CCE} & Uplink scheduling failure rate caused by CCE allocation failure & MAC \\
{UL\_SCHED\_FAIL\_CNT\_CCE} & Count of uplink scheduling failures due to insufficient CCE resources & MAC \\
{DL\_SCHED\_FAIL\_RATE\_CCE} & Downlink scheduling failure rate caused by CCE allocation failure & MAC \\
{DL\_SCHED\_FAIL\_CNT\_CCE} & Count of downlink scheduling failures due to CCE shortage & MAC \\
{DL\_CCE\_FAIL\_CNT\_TOTAL} & Total count of downlink CCE allocation failures caused by overall CCE shortage & MAC \\
{DL\_CCE\_FAIL\_RATE\_TOTAL} & Failure rate of downlink CCE allocation due to overall CCE shortage & MAC \\
{DL\_CCE\_FAIL\_CNT\_QUOTA} & Count of downlink CCE allocation failures due to exceeding quota limits & MAC \\
{DL\_CCE\_FAIL\_CNT\_CONFLICT} & Count of downlink CCE allocation failures due to search space conflicts & MAC \\
{UL\_CCE\_FAIL\_CNT\_TOTAL} & Total count of uplink CCE allocation failures caused by overall CCE shortage & MAC \\
{UL\_CCE\_FAIL\_RATE\_TOTAL} & Failure rate of uplink CCE allocation due to overall CCE shortage & MAC \\
{UL\_CCE\_FAIL\_CNT\_CONFLICT} & Count of uplink CCE allocation failures caused by search space conflicts & MAC \\
{DL\_RLC\_TPUT} & Downlink Radio Link Control (RLC) layer throughput in Mbit & RLC \\
{DL\_RLC\_LASTTTI\_RATIO} & Ratio of downlink RLC traffic in the last Transmission Time Interval (TTI) & RLC \\
{UL\_RLC\_TPUT} & Uplink RLC layer throughput in Mbit & RLC \\
{UL\_RLC\_SMALLPKT\_RATIO} & Ratio of small packet traffic in the uplink RLC layer & RLC \\
{RLC\_UL\_LATENCY} & Average latency of uplink RLC packets (ms) & RLC \\
{RLC\_DL\_LATENCY} & Average latency of downlink RLC packets (ms) & RLC \\
{UL\_RLC\_RETX\_SDU} & Number of retransmitted SDU segments in uplink RLC & RLC \\
{UL\_RLC\_SDU} & Total number of transmitted SDUs in uplink RLC & RLC \\
{PDCP\_DL\_TPUT} & Total downlink PDCP layer throughput (Mbit) & PDCP \\
{PDCP\_DL\_LATENCY} & Average latency of downlink PDCP packets (ms) & PDCP \\
{PDCP\_UL\_TPUT} & Total uplink PDCP layer throughput (Mbit) & PDCP \\
{PDCP\_UL\_LATENCY} & Average latency of uplink PDCP packets (ms) & PDCP \\
\hline
\end{tabular}
\end{table*}
\FloatBarrier

\subsection{Representative Rule-Based Labeling Example}

The practical rule-based annotator consists of category-specific logic defined on a subset of KPIs. To make this process concrete, we present the rule for \textit{Uplink Weak Coverage} as a representative example. 

A sample is assigned the label \textit{Uplink Weak Coverage} if it satisfies all of the following conditions:
\begin{itemize}
\item[(1)] \texttt{PDCP\_DL\_LATENCY} $\ge$ 200 ms or \texttt{PDCP\_UL\_LATENCY} $\ge$ 200 ms;
\item[(2)] \texttt{RLC\_DL\_LATENCY} $\ge$ 200 ms or \texttt{RLC\_UL\_LATENCY} $\ge$ 200 ms;
\item[(3)] 
\[
\frac{\texttt{UL\_RLC\_RETX\_SDU}}{\texttt{UL\_RLC\_SDU} + \epsilon} > 0.1, \quad \epsilon = 10^{-5};
\]
\item[(4)] At least one of:
\begin{itemize}
  \item[] \texttt{DL\_RBLER} $\ge$ 0.1,
  \item[] \texttt{UL\_RBLER} $\ge$ 0.1,
  \item[] \texttt{UL\_DTX\_Ratio} $\ge$ 0.2,
  \item[] \texttt{UL\_MAC\_HARQ\_RETX\_MAX} $\ge$ 50, 
  \item[] \texttt{DL\_MAC\_HARQ\_RETX\_MAX} $\ge$ 50;
\end{itemize}
\item[(5)] \texttt{UL\_DMRS\_RSRP\_MIN} $\le$ -125 dBm or \texttt{UL\_SRS\_RSRP} $\le$ -130 dBm.
\end{itemize}


\subsection{Evaluation Metrics}
\label{appendix:metrics}

Let $N$ denote the total number of test samples, $\hat{y}_i$ the predicted label and $y_i$ the ground-truth label for the $i$-th sample, and $C = 6$ the number of root-cause categories. For class $c$, let $\mathrm{TP}_c$, $\mathrm{FP}_c$, and $\mathrm{FN}_c$ denote true positives, false positives, and false negatives, respectively.

\noindent\textbf{Classification Accuracy} measures the proportion of correctly classified samples:
\begin{equation}
\text{Accuracy} = \frac{\sum_{i=1}^{N} \mathbb{1}(\hat{y}_i = y_i)}{N},
\end{equation}
where $\mathbb{1}(\cdot)$ is the indicator function.

\noindent\textbf{Macro-averaged F1-score} averages the per-class F1 scores equally across all $C$ categories, giving equal weight to each class regardless of support:
\begin{equation}
\text{Macro-F1} = \frac{1}{C}\sum_{c=1}^{C} F1_c.
\end{equation}

\noindent\textbf{Confusion Matrix} provides a class-level view of prediction outcomes. Let $M \in \mathbb{R}^{C \times C}$ denote the confusion matrix, where $M_{ij}$ is the number of samples whose true label is class $i$ and predicted label is class $j$. The diagonal entries represent correct identifications, whereas the off-diagonal entries reveal confusion between different root causes.


\section{Baseline Settings}
\label{appendix:baselines}

\noindent\textit{Classical ML Methods:}
\begin{itemize}
    \item \textbf{Sliding Window k-Nearest Neighbors (Sliding-KNN):} A non-parametric method based on distance metrics applied to fixed-length subsequences of the time series.
    \item \textbf{Support Vector Machine (SVM):} A classical supervised learning algorithm using the RBF kernel.
    \item \textbf{XGBoost:} A gradient-boosted decision tree ensemble that performs feature-level splits on the flattened time-series representation.
    \item \textbf{Time Series Forest (TSF):} An ensemble learning method that trains multiple decision trees on randomly selected intervals of the input time series.
    \item \textbf{Bag-of-Patterns (BOP):} A symbolic representation-based method that first discretizes time series into symbolic words and then applies histogram-based feature extraction for classification.
\end{itemize}
For all ML baselines, the time-series data are reshaped into two-dimensional arrays of shape $(n_\text{samples}, n_\text{timesteps} \times n_\text{features})$ before feeding into the models.

\noindent\textit{Semi-supervised Learning Methods:}
\begin{itemize}
    \item \textbf{TS-TCC~\cite{eldele2021time}:} A semi-supervised framework consisting of unsupervised contrastive learning with traditional time-series augmentations, followed by fine-tuning using a small amount of labeled data.
    \item \textbf{CA-TCC~\cite{eldele2023self}:} Extends TS-TCC with a four-stage semi-supervised training pipeline that generates pseudo labels and then refines the model through supervised contrastive learning.
\end{itemize}


{\color{R3color}
\noindent\textit{Weakly-Supervised Learning Methods:}
\begin{itemize}
    \item \textbf{Co-teaching~\cite{han2018coteaching}:} Two networks are trained simultaneously, and each selects small-loss samples to teach the other, thereby reducing the impact of noisy rule-based labels.
    \item \textbf{DivideMix~\cite{li2020dividemix}:} A representative noisy-label learning method that uses a Gaussian mixture model to split samples into clean and noisy sets, and then trains them in a semi-supervised manner with two networks.
    \item \textbf{Scale-teaching~\cite{liu2023scale}:} A multi-scale time-series noisy-label learning method that learns discriminative patterns across different temporal scales, performs cross-teaching for small-loss sample selection, and corrects large-loss samples via graph-based label propagation.
\end{itemize}
}

% \noindent\textit{Supervision settings:} (1) Classical ML methods are trained on expert labels only. (2) TS-TCC and CA-TCC treat rule-labeled samples as unlabeled data and fine-tune on expert labels. (3) CNN-full is supervised by both label sources. (4) {\color{R3color}Weakly supervised baselines use rule labels with their respective noise-handling mechanisms.} For encoder-based DL methods, all models share the same classification head (a single fully connected linear layer).


\section{Diffusion Validation Details}

\subsection{Augmentation Baselines and Evaluation Metrics for Diffusion Validation}
\label{appendix:diffusion_baselines}


\noindent\textit{Augmentation Baselines:}
\begin{itemize}
    \item \textbf{DK-Root (Proposed):} For the proposed conditional diffusion-based augmentation, the weak and strong augmentation time steps are sampled as:
    \begin{align}
        t_{\text{weak}} &\sim \mathcal{U}\left(0,\ \left\lfloor \frac{T}{5} \right\rfloor \right) \\
        t_{\text{strong}} &\sim \mathcal{U}\left(\left\lfloor \frac{T}{5} \right\rfloor,\ \left\lfloor \frac{T}{2} \right\rfloor \right)
    \end{align}

    {\color{R2color}
    \item \textbf{TimeGAN~\cite{yoon2019timegan}:} A GAN-based time-series generation framework that jointly learns an embedding space and a generative adversarial model with supervised temporal dynamics. Class labels are injected via a learnable embedding layer concatenated with the latent noise input to the generator. Weak and strong augmentations are controlled via the noise temperature $\tau$:
    \[
        \mathbf{z}_{\text{aug}} = \tau \cdot \mathbf{z}, \quad \mathbf{z} \sim \mathcal{N}(0, \mathbf{I})
    \]
    where $\tau_{\text{weak}} = 0.6$ and $\tau_{\text{strong}} = 1.2$.
    }

    \item \textbf{Noise Injection:} Gaussian noise scaled by per-channel standard deviation:
    \[
        \tilde{x} = x + 0.1 \cdot \sigma \cdot \mathcal{N}(0, 1)
    \]

    \item \textbf{Scaling:} Channel-wise multiplicative perturbation:
    \[
        \tilde{x}_{i,t} = x_{i,t} \cdot \alpha_t, \quad \alpha_t \sim \mathcal{N}(1, \sigma^2),\ \sigma = 1.1
    \]
\end{itemize}

\noindent\textit{Evaluation Metrics:}
\begin{itemize}
    \item \textbf{Average Inter-class Distance}: Mean Euclidean distance between class centers in feature space. Larger values indicate better separability.
    \item \textbf{Silhouette Score~\cite{rousseeuw1987silhouettes}}: Measures intra-cluster cohesion vs.\ inter-cluster separation. Ranges from $-1$ to $1$; higher is better.
    \item \textbf{Calinski-Harabasz Index (CH)~\cite{calinski1974dendrite}}: Ratio of between-cluster to within-cluster dispersion. Higher values indicate more compact and well-separated clusters.
    \item \textbf{Davies-Bouldin Index (DB)~\cite{davies2009cluster}}: Average similarity of each cluster to its nearest cluster. Lower values imply better inter-cluster distinction.
    \item \textbf{Mutual Information (MI)}: Measures dependency between features and class labels. Higher MI indicates stronger semantic alignment with ground-truth labels.
    \item \textbf{Intra-class Variance}: Average variance of feature vectors within each class. Lower values indicate higher intra-class compactness and semantic consistency.
\end{itemize}

}


}

% \clearpage
\bibliographystyle{IEEEtran}
\bibliography{ref}

\end{document}